% discussion_encodings.tex
\subsection{Methodological Implications}

The clearest lesson concerns reverse keying, the Introduction's fault
line: a reverse-scored response keeps its content on an inverted scale; an
embedding multiplied by $-1$ is an anti-topic vector. Hence
\texttt{atomic\_reversed} matched keyed human correlations worst of the
cosine family ($r = .137$ vs.\ .434 unflipped), its checkerboard signs
driving split-half reliability negative and rendering disattenuated
congruence undefined---a pathology-diagnosing NA, reported, not
suppressed. This substantiates the caution that embeddings do not
differentiate reverse-keyed items \citep{suarez-alvarez-etal-2026},
explains why keyed sign-flipping fared worst in HEXACO pseudo-factor
analyses \citep{guenole-etal-2025}, fits the positive-sign bias surviving
polarity fine-tuning \citep{hommel-arslan-2025}, and settles negatively
whether an embedding can be meaningfully ``reversed''
\citep{pellert-etal-2026}. Content-not-key logic recurs package-wide:
anchoring and projection both flagged N12 (``I am relaxed most of the
time.''), reverse-keyed yet semantically the inventory's most calm-pole
statement (.82). Three upshots follow: prefer keying-free encodings
(mean-centered Pearson, SQuID) or signed NLI over sign-flipping; build
candidate-vetting reference fits without the scoring column, keeping
centroids purely topical; and drop earlier short-form pipelines'
drift-prone manual paraphrasing of reverse-keyed items
\citep{jung-seo-2025,wang-etal-2026}, checking substitutions scale by
scale.

Second, no single encoding or congruence metric dominates; each
operationalizes a distinct inter-item relation---cosine symmetric and
graded \citep{reimers-gurevych-2019}, entailment directional over ordered
premise--hypothesis pairs \citep{bowman-etal-2015}. Profiles mirror
accordingly: signed NLI best matched the theoretical partition
(NMI $= .548$) and raw correlations ($r = .561$) but worst matched loading
shape ($\bar{\phi} = .310$ vs.\ cosine's .744), while SQuID won keyed
pair-level agreement ($r = .567$) on a near-unfactorable differenced
matrix (KMO $= .022$). What language predicts is thus relation- and
encoding-dependent, echoing item-parameter modeling: difficulty is
text-predictable \citep{peters-etal-2025}, discrimination largely not
\citep{han-etal-2025}, though simulation-based item-characteristic-curve
reconstruction recovers it more readily \citep{ormerod-2026}. Two
reporting practices follow: compare like with like---keying-free encodings
to raw correlations, keyed to rescored; conflation under- or overstates
coherence (NLI: $r_{\text{raw}} = .561$ vs.\
$r_{\text{keyed}} = .063$)---and report congruence on several metrics
(loading shape, partition agreement, pair-level correlation), which can
rank encodings oppositely.

Third, semantic-space fit needs matched chance baselines: by survey
conventions the RMSR of .032 reads as good yet exceeded the entire
structureless-embedding null range (.0119--.0134), the CAF (.412) fell
below its nulls (.501--.532), and the TEFI ($-27.71$) beat all 100 null
replicates ($-21.66$ to $-17.82$)---two of three conventional readings
invert. This generalizes beyond confirmatory modeling the finding that
chance meets conventional thresholds alarmingly often in embedding space
\citep{pokropek-2026} and answers the ``no sample size'' problem of
\citet{suarez-alvarez-etal-2026}: absent an $N$, Monte Carlo nulls supply
empirical reference distributions; we propose such calibrated reporting as
standard for semantic factor analyses. Retention repeats the pattern:
semantic criteria disagree (parallel 5, Kaiser 6, TEFI 2, EGA 5) as
simulations predict \citep{garrido-etal-2025, golino-2026}, while
human-side parallel analysis at $N = 874{,}434$ endorsed ten factors as
trivially small minor factors cleared the noise threshold; such two-sided
minor-factor proliferation is precisely the case for consensus-based
retention with theory-fixed solutions \citep{goldberg-1990}.
